\addchap{Preface and Notes}
\label{ch:Preface}

\subsection*{Use of External Tools}
For the linguistic drafting and correction of the Introduction (\cch{Einführung}) and the Execution (\cch{Durchfuehrung}), the AI \textit{Lumo AI} as well as \textit{Duden Mentor} were employed. However, the substantive evaluation and the actual creation of the document were carried out entirely without the assistance of Artificial Intelligence, unless explicitly marked otherwise. \cite{Duden,Lumo}

All graphics contained in the document that lack a specific reference were independently created or edited in Inkscape, or are plots generated in Python. Both the Python code and the LaTeX document were prepared in the Visual Studio Code (VS Code) development environment.

% In the area of code generation, AI was used supportively. Should a code cell be primarily generated by an AI, this was explicitly noted within the respective cell to ensure transparency.

\subsection*{Own Code and Transparency}
To carry out the evaluations, a custom Python library was developed. The entire source code - both the Python scripts and the LaTeX sources - is self-created and publicly available online.

The entire project structure can be found on my GitHub. Here, both the Python code and the LaTeX document code are documented. \cite{githubPAP2}

The Python functions used are specifically identifiable on GitHub under the name "Papulator" \cite{Papulator}

\subsection*{Use of University Materials}
This document is based on the script by Jens Wagner (as of April 27, 2026). Graphics from this script were adopted for the Introduction. These excerpts are transparently marked; the figures used always originate from the chapters of the script that share the same experiment number and experiment name as in the present document. In this experiment, the graphics from \emph{Experiment \versuchsnummer \versuchsname{\versuchsnummer}} were used. \cite{skriptPAP21}